%!TeX root=MemoriaTFG.tex

\chapter*{Resumen}
\addcontentsline{toc}{chapter}{Resumen}
\markboth{RESUMEN}{RESUMEN}
\vspace*{-1.5cm}

Este trabajo evalúa el rendimiento y la capacidad de generalización
de los modelos fundacionales dermatológicos disponibles, en función
del volumen de datos etiquetados, del tipo de supervisión, de la
heterogeneidad taxonómica entre conjuntos de datos y de la
disparidad entre fototipos cutáneos. Su respuesta condiciona la
viabilidad de integrarlos en teledermatología hospitalaria con
cohortes mediterráneas en español, escenario donde las referencias
recientes del grupo de la Universidad Monash ---PanDerm, la familia
DermLIP sobre Derm1M y DermFM-Zero--- carecen de evaluación
independiente comparable a la fecha de cierre.

Para responderla, se evalúan catorce modelos fundacionales sobre
doce conjuntos de datos externos públicos, armonizados mediante una
ontología jerárquica de tres niveles ($4$ categorías etiológicas,
$43$ subcategorías clínicas y $367$ diagnósticos) codificada en
SNOMED~CT y CIE-10. La evaluación combina cuatro protocolos
---sondeo lineal, ajuste fino supervisado, segmentación binaria
promptable y clasificación zero-shot multimodal--- ejecutados sobre
una NVIDIA DGX Spark (chip Grace Hopper GB10, $128$\,GB unificados,
BF16), plataforma aarch64 no habitual en la literatura.

PanDerm Large se confirma como el codificador congelado más fiable,
con una mejora media de $+9{,}0$\,pp en exactitud balanceada
respecto a PanDerm Base, y satura con un volumen reducido de datos
etiquetados: el $1\%$ del conjunto de entrenamiento de HAM10000
recupera el $96{,}2\%$ del AUROC obtenido con el conjunto completo.
La evaluación zero-shot multimodal queda $29{,}4$\,pp por debajo del
sondeo lineal supervisado en exactitud balanceada de nivel~L1 sobre
el dataset DermapixelAI~1.0 (PanDerm Large $0{,}735$ frente a
$0{,}441$ de la mejor configuración zero-shot), lo que acota el
coste de prescindir de adaptación supervisada. La auditoría por hash
MD5 detecta que uno de los doce datasets de evaluación, Dermnet,
está íntegramente contenido ($100\%$) en el preentrenamiento de
Derm1M, por lo que sus cifras zero-shot previas deben interpretarse
con \emph{caveat}. La validación sobre DermapixelAI~1.0 confirma la
transferencia a material clínico no anglosajón sin reentrenar la
rama lingüística. Leído en conjunto, el estudio no identifica un
único modelo óptimo, sino una especialización por tarea y nivel
ontológico ---la ventaja de PanDerm Large no se extiende
uniformemente a la subcategoría clínica L2 ni a la cola larga
L3---, lo que desplaza el cuello de botella del avance desde el
codificador hacia la taxonomía, las etiquetas, el alineamiento
texto-imagen, los \emph{prompts} y la equidad.

El trabajo aporta tres contribuciones reproducibles principales: un
benchmark comparativo homogéneo de catorce modelos fundacionales
bajo un mismo protocolo, el dataset DermapixelAI~1.0 ---$1\,089$ imágenes y $669$
casos clínicos con texto narrativo en castellano, bajo licencia
CC-BY-NC-SA~4.0--- y la auditoría de solapamiento sobre los doce
datasets. Sobre esta base se introduce Dermapixel R0, el
primer modelo fundacional dermatológico en español, inicializado
desde PanDerm y adaptado con DermapixelAI~1.0 como versión base; su
escalado completo queda como línea futura.

\vspace*{-0.4cm}
\section*{Palabras clave}
modelos fundacionales, dermatología clínica, PanDerm, DermLIP,
evaluación empírica, ontología, equidad, heterogeneidad taxonómica.

% ----------------------------------------------------------------------

\chapter*{Abstract}
\addcontentsline{toc}{chapter}{Abstract}
\markboth{ABSTRACT}{ABSTRACT}
\vspace*{-1.5cm}

This work evaluates the performance and generalisation capacity of
the available dermatology foundation models, as a function of the
labelled data budget, the supervision regime, taxonomic
heterogeneity across datasets and skin phototype disparity. Its
answer conditions the
viability of integrating these models into hospital teledermatology
workflows on Spanish-speaking Mediterranean cohorts, a scenario in
which the recent references from the Monash group ---PanDerm, the
DermLIP family over Derm1M and DermFM-Zero--- lack an independent
comparable evaluation at the time of writing.

To answer it, fourteen foundation models are evaluated over twelve
public external datasets, harmonised through a three-level
hierarchical ontology ($4$ aetiological families, $43$ clinical
subcategories and $367$ diagnoses) encoded in SNOMED~CT and ICD-10.
The evaluation combines four protocols ---linear probing, supervised
fine-tuning, promptable binary lesion segmentation and multimodal
zero-shot classification--- executed on a NVIDIA DGX Spark with a
Grace Hopper GB10 chip ($128$\,GB unified, BF16 precision), an
aarch64 platform uncommon in the experimental literature.

PanDerm Large stands out as the most reliable frozen encoder, with a
mean improvement of $+9{,}0$\,pp in balanced accuracy over PanDerm
Base, and saturates with a small labelled budget: $1\%$ of the
HAM10000 training set recovers $96{,}2\%$ of the AUROC obtained with
the full set. Multimodal zero-shot evaluation falls $29{,}4$\,pp
below supervised linear probing in L1 balanced accuracy on the
DermapixelAI~1.0 dataset (PanDerm Large $0{,}735$ versus $0{,}441$
for the best zero-shot configuration), bounding the cost of forgoing
supervised adaptation. The MD5-based overlap audit detects that one
of the twelve evaluation datasets, Dermnet, is fully contained
($100\%$) in the Derm1M pretraining set, so its prior zero-shot
figures must be read with this \emph{caveat}. Validation on
DermapixelAI~1.0 confirms model transfer to non-Anglophone clinical
material without retraining the language branch. Taken together, the
study does not identify a single optimal model but a specialisation
by task and ontological level ---PanDerm Large's advantage does not
extend uniformly to the intermediate L2 subcategory or the L3 long
tail---, shifting the bottleneck of progress from the encoder towards
taxonomy, labels, text-image alignment, \emph{prompts} and fairness.

The work contributes three reproducible items: a homogeneous
comparative benchmark of the fourteen families under a common
protocol, the DermapixelAI~1.0 dataset ---$1\,089$ images linked to
$669$ clinical cases with narrative text and synthetic captions in
Spanish, under a CC-BY-NC-SA~4.0 licence--- and the overlap audit
across the twelve evaluation sets. Building on this basis,
Dermapixel R0 is introduced: a Spanish dermatology foundation model
initialised from PanDerm and subsequently adapted using the
DermapixelAI~1.0 dataset as an initial base, whose full scaling is
left as future work.

\vspace*{-0.4cm}
\section*{Keywords}
foundation models, clinical dermatology, PanDerm, DermLIP, empirical
evaluation, ontology, fairness, taxonomic heterogeneity.
